% ============================================================
\section{Reinforcement Learning Training and Evaluation}

This section documents \texttt{05\_train\_rl.py} (training) and \texttt{evaluate.py} (evaluation). Together they implement a policy-gradient method that trains a Graph Neural Network (GNN) to make variable-selection decisions inside SCIP's branch-and-bound solver, with the objective of minimising the total number of nodes explored.

% ============================================================
\subsection{Branch-and-Bound as a Markov Decision Process}

Branch-and-bound solves a MIP by recursively splitting the feasible region along integer variables. At each \emph{branching node} the solver must choose a fractional variable to branch on --- this is the \emph{variable selection} (or \emph{branching}) decision. The quality of this choice directly determines how quickly the tree can be pruned and how many nodes must be processed before the optimal solution is certified.

The paper casts variable selection as an MDP with the following components:

\begin{description}
    \item[State $s_t$] The \texttt{NodeBipartite} observation produced by Ecole at the current open node: a bipartite graph whose left nodes are the active constraints (with LP-derived features) and whose right nodes are the candidate variables (with LP, branching-history, and objective features). Edge features encode the constraint matrix coefficients.
    \item[Action $a_t$] An integer index into \texttt{action\_set}, the set of LP-fractional candidate variables at the current node.
    \item[Reward $r_t$] A scalar that depends on the training mode (see Section~\ref{sec:training_modes}); in all modes it is proportional to the \emph{negative} number of nodes explored, so maximising the expected return corresponds to finding a smaller search tree.
    \item[Episode] One full branch-and-bound solve of a single MIP instance. An episode ends when SCIP certifies optimality (or exceeds the time limit).
\end{description}

% ============================================================
\subsection{Policy: the GNN Actor}

The policy $\pi_\theta$ is a Graph Neural Network (\texttt{GNNPolicy} in \texttt{actor/actor.py}). It takes the \texttt{NodeBipartite} bipartite graph as input and outputs a scalar logit for every candidate variable:

\begin{lstlisting}
logits = self.actor(
    batch.constraint_features,   # constraint node features
    batch.edge_index,            # bipartite adjacency
    batch.edge_attr,             # edge (coefficient) features
    batch.variable_features      # variable node features
)
logits = logits[batch.action_set]   # restrict to LP candidates
\end{lstlisting}

Action selection from these logits is either:
\begin{itemize}
    \item \textbf{Greedy} (validation/evaluation): $a = \arg\max_i \, \texttt{logits}_i$.
    \item \textbf{Stochastic} (training): $a \sim \mathrm{Categorical}(\texttt{logits})$, i.e.\ a sample from the softmax distribution.
\end{itemize}

% ============================================================
\subsection{Training Modes}
\label{sec:training_modes}

Three training modes are implemented, each with a different MDP formulation. They share the same GNN architecture and REINFORCE update rule but differ in how episodes are truncated and how returns are computed.

\subsubsection{\texttt{mdp} --- Standard MDP}

The standard full-tree MDP. The environment (\texttt{MDPBranchingEnv}) runs SCIP with its default node selector. The return assigned to a transition collected at step $t$ (node $n_t$) is:

\[
    G_t = c_{t} - c_{\text{final}}
\]

where $c_t$ is the cumulative node count at step $t$ and $c_{\text{final}}$ is the cumulative node count at episode end. Because $c_t \leq c_{\text{final}}$, this is always $\leq 0$, and equals zero only if no further nodes are opened after branching at $n_t$. Equivalently, $G_t$ equals minus the number of nodes opened \emph{after} step $t$.

\begin{lstlisting}
# credit assignment for mdp
for transition in transitions:
    transition.returns = transition.cum_nnodes - cum_nnodes
\end{lstlisting}

\subsubsection{\texttt{tmdp+DFS} --- Truncated MDP with Depth-First Search}

During training the node selector is overridden to Depth-First Search (DFS), which forces SCIP to explore one branch completely before backtracking (\texttt{DFSBranchingDynamics}). This truncates the episode to a single root-to-leaf path, making trajectories shorter and variance lower. The return is the \emph{negative subtree size} rooted at the branching node:

\[
    G_t = -|\mathrm{subtree}(n_t)| - 1
\]

The \texttt{TreeRecorder} class computes subtree sizes in a bottom-up pass over the recorded branching decisions. A DFS trajectory visits every node in the explored subtree exactly once, so this return directly measures the cost of the branching decision.

\begin{lstlisting}
# credit assignment for tmdp+DFS and tmdp+ObjLim
subtree_sizes = tree_recorder.calculate_subtree_sizes()
for transition in transitions:
    transition.returns = -subtree_sizes[transition.node_id] - 1
\end{lstlisting}

\subsubsection{\texttt{tmdp+ObjLim} --- Truncated MDP with Objective Limit}

Instead of changing the node selector, this mode injects a \emph{primal bound} (objective limit) into SCIP at the start of each episode (\texttt{ObjLimBranchingDynamics}). The bound is set to the pre-computed optimal value $z^*$ plus a small slack $\varepsilon$:

\begin{lstlisting}
# in train_batch_generator()
eps = -0.1 if maximization else 0.1
yield [{'path': instance,
        'sol': train_sols[instance] + eps,   # primal bound
        'seed': ...}
       for instance in ...]
\end{lstlisting}

Any node whose LP relaxation cannot improve beyond this bound is pruned immediately by SCIP. This truncates the tree from below: only nodes that could potentially discover a solution better than $z^* + \varepsilon$ are explored, which typically reduces the tree to the left-most branch plus a handful of pruned siblings. Credit assignment is the same subtree-size formula as \texttt{tmdp+DFS}.

% ============================================================
\subsection{Concurrent Agent Architecture}

\subsubsection{Overview of Threads}

Running one Ecole environment at a time would leave the GPU idle while SCIP is branching. The code uses a multi-threaded design with three kinds of threads (see \texttt{agent.py}):

\begin{enumerate}
    \item \textbf{Agent threads} (\texttt{num\_agents = 12}): Each agent owns one Ecole branching environment. At every branching node it extracts the state, sends a policy query to the \texttt{PolicySampler}, waits for the action, then steps the environment.
    \item \textbf{PolicySampler thread}: Drains the \texttt{policy\_queries\_queue}, accumulates all waiting requests into a batch, calls \texttt{brain.sample\_action\_idx()} (a single GPU forward pass), and distributes the resulting action indices back to the requesting agents.
    \item \textbf{Main thread}: Manages the epoch loop, calls \texttt{brain.update()}, triggers validation, and saves checkpoints.
\end{enumerate}

The key benefit is \emph{GPU batching}: because all 12 agents submit policy queries concurrently, the PolicySampler can process them in one batch rather than 12 serial forward passes.

\subsubsection{Preemptive Job Scheduling}

Jobs are submitted with \texttt{block\_policy=True}, which creates a \texttt{threading.Event} (\texttt{policy\_access}) that is initially \emph{unset}. Agents load their instances and run SCIP until the first branching node, then block on \texttt{policy\_access.wait()}. The main thread sets \texttt{policy\_access.set()} only when it is ready to service the epoch, at which point all blocked agents proceed simultaneously. This overlaps SCIP's instance-loading time with the neural-network update of the previous epoch:

\begin{lstlisting}
# start next epoch's jobs before processing current epoch's results
if is_training_epoch(epoch + 1):
    t_samples_next, ..., t_access_next = \
        agent_pool.start_job(train_batch, ..., block_policy=True)

# ... process and update current epoch ...

# release the preempted agents
t_access.set()
\end{lstlisting}

% ============================================================
\subsection{REINFORCE Policy Gradient Update}

After all training episodes in an epoch finish, the main thread calls \texttt{brain.update(transitions)}. The loss combines a REINFORCE term and an entropy bonus:

\[
    \mathcal{L}(\theta) = -\frac{1}{N} \sum_{i=1}^{N} G_i \cdot \log \pi_\theta(a_i \mid s_i)
                          \;-\; \beta \cdot \frac{1}{N}\sum_{i=1}^{N} H[\pi_\theta(\cdot \mid s_i)]
\]

where $N$ is the total number of collected transitions, $G_i$ is the return, and $H$ denotes the entropy of the action distribution. The entropy bonus (weight $\beta = 10^{-5}$) prevents premature collapse to a deterministic policy during training. The optimizer is Adam with learning rate $10^{-6}$.

\begin{lstlisting}
# REINFORCE term
reinforce_loss = - (returns * dist.log_prob(batch.action)).sum() / n_samples
# Entropy bonus (subtracted to maximise entropy)
entropy = dist.entropy().sum() / n_samples
loss = reinforce_loss - entropy_bonus * entropy
loss.backward()
optimizer.step()
\end{lstlisting}

Only a fraction \texttt{sample\_rate = 0.05} of branching decisions is collected per episode to reduce memory usage and decorrelate samples.

% ============================================================
\subsection{Training Loop}

\begin{enumerate}
    \item Read \texttt{config.default.json} and apply any command-line overrides.
    \item Load training instances and pre-computed optimal solutions (\texttt{instance\_solutions.json}).
    \item Instantiate \texttt{Brain} (GNNPolicy + Adam optimizer) and \texttt{AgentPool} (12 agent threads + PolicySampler).
    \item For each epoch $e = 0, \ldots, T$ (default $T = 20000$):
    \begin{enumerate}
        \item \textbf{Validation} (every 50 epochs): run greedy policy on 20 instances $\times$ 5 seeds. Compute 1-shifted geometric mean of nnodes. Save model if this is the best value so far.
        \item \textbf{Training}: sample a random batch of 10 training instances; collect REINFORCE transitions; call \texttt{brain.update()}.
    \end{enumerate}
    \item Stop after 6 days of wall-clock time regardless of epoch count.
\end{enumerate}

The model checkpoint is written to \texttt{actor/<problem>/0/<timestamp>--best\_params--<mode>.pkl}.

% ============================================================
\subsection{Evaluation (\texttt{evaluate.py})}

\subsubsection{Benchmark Design}

Evaluation tests every combination of policy $\times$ instance $\times$ seed, writing all results to a CSV file. Two instance sets are used:

\begin{itemize}
    \item \textbf{Test}: 40 instances drawn from the same distribution as training.
    \item \textbf{Transfer}: 40 instances from a harder distribution (e.g.\ larger graphs or more knapsacks), testing out-of-distribution generalisation.
\end{itemize}

Each instance is solved with 5 random seeds (0--4) to account for SCIP's internal stochasticity.

\subsubsection{Policies Evaluated}

\paragraph{SCIP internal brancher (\texttt{relpscost}).}
Reliability Pseudocost Branching is SCIP's state-of-the-art default rule. It is run via the \texttt{Configuring} environment: the brancher priority is raised to the highest level, then \texttt{env.step(\{\})} triggers a full solve in one call (see Section~3 for how this works internally).

\paragraph{GCNN models (\texttt{il}, \texttt{mdp}, \texttt{tmdp+DFS}, \texttt{tmdp+ObjLim}).}
Each GNN model is run via the \texttt{Branching} environment. At every branching node the agent observes \texttt{NodeBipartite} features and selects the variable with the highest logit (greedy):

\begin{lstlisting}
observation, action_set, _, done, _ = env.reset(instance['path'])
while not done:
    # observation -> tensor features
    logits = policy['model'](*observation)
    action = action_set[logits[action_set.astype(np.int64)].argmax()]
    observation, action_set, _, done, _ = env.step(action)
\end{lstlisting}

\subsubsection{Metrics Recorded}

For each instance-policy-seed triple the following metrics are written to the CSV:

\begin{center}
\begin{tabular}{ll}
\toprule
\textbf{Field} & \textbf{Meaning} \\
\midrule
\texttt{nnodes}   & Number of branch-and-bound nodes processed \\
\texttt{nlps}     & Number of LP relaxations solved \\
\texttt{stime}    & SCIP solving time (s) \\
\texttt{gap}      & Optimality gap at termination \\
\texttt{status}   & SCIP termination status (optimal, timelimit, \ldots) \\
\texttt{walltime} & Wall-clock time including Python overhead \\
\texttt{proctime} & CPU process time \\
\bottomrule
\end{tabular}
\end{center}

The primary metric for comparing policies is \texttt{nnodes}: fewer nodes means the learned branching policy pruned the tree more efficiently than the baseline. The time limit is 3600 seconds per instance; instances not solved to optimality within this budget are recorded with their current gap.
